
\documentclass{article}
\usepackage{colortbl}
\usepackage{makecell}
\usepackage{multirow}
\usepackage{supertabular}

\begin{document}

\newcounter{utterance}

\twocolumn

{ \footnotesize  \setcounter{utterance}{1}
\setlength{\tabcolsep}{0pt}
\begin{supertabular}{c@{$\;$}|p{.15\linewidth}@{}p{.15\linewidth}p{.15\linewidth}p{.15\linewidth}p{.15\linewidth}p{.15\linewidth}}

    \# & $\;$A & \multicolumn{4}{c}{Game Master} & $\;\:$B\\
    \hline 

    \theutterance \stepcounter{utterance}  

    & & \multicolumn{4}{p{0.6\linewidth}}{\cellcolor[rgb]{0.9,0.9,0.9}{%
	\makecell[{{p{\linewidth}}}]{% 
	  \tt {\tiny [A$\langle$GM]}  
	 Eres un experto en juegos de puzles de física 2D. El entorno del juego es el siguiente:\\ \tt \\ \tt Hay una bola azul que se puede dejar caer desde el borde superior de la pantalla en cualquier posición horizontal (coordenada x). La posición de caída vertical siempre está fija en y = 100.\\ \tt La bola cae por gravedad y puede rodar o rebotar en los obstáculos colocados a lo largo del nivel.\\ \tt El objetivo es que la bola azul alcance y toque una bola objetivo verde ubicada entre los obstáculos.\\ \tt Los obstáculos pueden ser paredes, rampas u otras formas físicas que afectan el movimiento de la bola.\\ \tt \\ \tt Esta es la configuración actual del juego:\\ \tt 1. La bola objetivo verde se encuentra en una posición fija conocida.\\ \tt 2. Las posiciones y formas de los obstáculos son fijas y permanecen invariables.\\ \tt \\ \tt Te proporciono una imagen como fotograma clave para este intento. La imagen incluye la siguiente información:\\ \tt 1. Los obstáculos presentes en la configuración del juego y sus posiciones.\\ \tt 2. La ubicación de la bola objetivo verde en el entorno del juego.\\ \tt \\ \tt Con base en esta imagen y la información anterior, sugiera la mejor coordenada x entera entre 100 y 500 donde debería soltarse la pelota azul la próxima vez para maximizar la probabilidad de alcanzar y tocar la pelota verde.\\ \tt IMPORTANTE: Responda solo con el formato: SOLTAR: <x>, donde <x> es un número entre 100 y 500. 
	  } 
	   } 
	   } 
	 & \\ 
 

    \theutterance \stepcounter{utterance}  

    & \multicolumn{4}{p{0.6\linewidth}}{\cellcolor[rgb]{0.8,1,0.9}{%
	\makecell[{{p{\linewidth}}}]{% 
	  \tt {\tiny [A$\rangle$GM]}  
	 SOLTAR: 300 
	  } 
	   } 
	   } 
	 & & \\ 
 

    \theutterance \stepcounter{utterance}  

    & & & \multicolumn{2}{p{0.3\linewidth}}{\cellcolor[rgb]{0.95,0.95,0.95}{%
	\makecell[{{p{\linewidth}}}]{% 
	  \tt {\tiny [GM$|$GM]}  
	 drop_x=300, collisions=11, goal=True, steps=377, dist=36.0 
	  } 
	   } 
	   } 
	 & & \\ 
 

\end{supertabular}
}

\end{document}
